\chapter{Cellulitis}

\section*{Definition}
Cellulitis is an acute, spreading bacterial infection of the deep dermis and subcutaneous tissue, presenting with erythema, warmth, edema, and tenderness. It lacks the well-demarcated, raised border of erysipelas, which is limited to the upper dermis and superficial lymphatics.

\section*{When to See a Doctor}
See your doctor if:
\begin{itemize}
  \item Rapidly expanding redness, severe or disproportionate pain, fever, or feeling very unwell
  \item Signs of systemic toxicity: tachycardia, hypotension, altered mental status
  \item Hemorrhagic bullae, skin necrosis, crepitus (suspect necrotizing fasciitis)
  \item Immunocompromised state, diabetes, or peripheral vascular disease
\end{itemize}

\section*{How It Develops}
Cellulitis typically follows disruption of the skin barrier, allowing bacteria to breach host defenses and proliferate in subcutaneous tissues. Bacterial toxins and cell wall components trigger a robust inflammatory response with widening of blood vessels, neutrophil chemotaxis, and immune signaling proteins release (IL-1, IL-6, TNF-\ensuremath{\alpha}), producing the classic signs of rubor, calor, tumor, and dolor. Lymphatic involvement may cause streaking and regional lymphadenopathy.

\section*{Causes}
\begin{itemize}
  \item \textbf{Bacterial:} \textit{Streptococcus pyogenes} (group A strep) is most common; \textit{Staphylococcus aureus} (including MRSA) in purulent cellulitis; Gram-negative rods in immunocompromised or diabetic people
  \item \textbf{Risk factors:} Venous insufficiency/lymphedema, diabetes mellitus, obesity, peripheral arterial disease, tinea pedis (interdigital entry), recent surgery or trauma
\end{itemize}

\section*{Related Symptoms}
\begin{itemize}
  \item Erysipelas, lymphangitis, abscess, necrotizing fasciitis, fever, leukocytosis
\end{itemize}
\textbf{Important:} Bilateral lower extremity erythema is rarely cellulitis; chronic venous stasis dermatitis is a more common mimic and should not be treated with antibiotics.

\section*{Medical Evaluation}
\begin{itemize}
  \item Clinical diagnosis based on history and physical examination; mark and date erythema borders
  \item Blood cultures if febrile, toxic, or immunocompromised; aspiration of leading edge for culture has low yield
  \item Imaging (CT or MRI) if necrotizing fasciitis or deep abscess is suspected
\end{itemize}

\section*{Treatment Options}
\begin{itemize}
  \item Antibiotic choice and duration depend on severity, allergy history, site, local resistance, and whether there is pus or an abscess; follow a clinician's prescription
  \item An abscess usually requires drainage; antibiotics may be added when there are systemic signs, extensive disease, or risk factors
  \item Severe infection or suspected necrotizing infection requires hospital assessment and intravenous treatment
  \item Treat underlying lymphedema with compression stockings after infection resolves
\end{itemize}

\section*{Self-Care and Home Management}
\begin{itemize}
  \item Elevate the affected limb to reduce edema
  \item Treat predisposing tinea pedis with antifungal cream
  \item Skin hygiene with gentle cleansing and emollients to prevent fissures
\end{itemize}

\section*{References}
\begin{thebibliography}{99}
\bibitem{cellulitis_infection:ref1} Stevens DL, Bisno AL, Chambers HF, et al. ``Practice Guidelines for the Diagnosis and Management of Skin and Soft Tissue Infections.'' Clin Infect Dis. 2014;59(2):147--159.
\bibitem{cellulitis_infection:ref2} Swartz MN. ``Cellulitis.'' N Engl J Med. 2004;350(9):904--912.
\bibitem{cellulitis_infection:ref3} Raff AB, Kroshinsky D. ``Cellulitis: A Review.'' JAMA. 2016;316(3):325--337.
\bibitem{cellulitis_infection:ref4} Gunderson CG, Chang JJ. ``Risk Factors for Recurrent Cellulitis: A Systematic Review and Meta-Analysis.'' J Am Acad Dermatol. 2021;85(3):665--675.
\bibitem{cellulitis_infection:ref5} Bailey E, Kroshinsky D. ``Cellulitis: Diagnosis and Management.'' Dermatol Ther. 2011;24(2):229--239.
\bibitem{cellulitis_infection:ref6} Brindle R, Williams OM, Barton E, et al. ``Assessment of Antibiotic Treatment for Cellulitis: A Systematic Review and Meta-Analysis.'' JAMA Dermatol. 2019;155(9):1033--1040.
\end{thebibliography}
